% FILE: 06_contribution_to_thesis.tex
% Project: lifecycle
% Role: process-lifecycle-state-definitions-and-transitions

\section{Contribution to the Dissertation}
\label{sec:lifecycle-contribution}

\subsection{Contribution to Prediction vs.~Coordination}
\label{sec:lifecycle-contribution-prediction}

The central thesis question is whether enterprise process intelligence is fundamentally
a prediction problem or a coordination problem. The lifecycle framework argues directly
for the coordination position. A process model at L1 (Ad-hoc) has no formal structure
and cannot support prediction: there is nothing to predict against. A process model at
L5 (Autonomic) does not predict future behavior in the statistical learning sense; it
coordinates present behavior through a receipted feedback loop that observes deviations,
analyzes their cause, plans a repair, and executes it---with every step producing a
typed artifact.

The distinction is not merely definitional. A prediction-based approach to process
governance asks: given historical data, what will happen next? A coordination-based
approach asks: given the declared process model and the observed event log, is the
current execution conforming, and if not, what action is required and by whom? The
lifecycle framework encodes the coordination approach in its gate criteria. Gate 3 does
not predict when fitness will fall below 0.95; it detects that it has fallen below 0.95
and triggers $T_{\mathrm{elastic}}$ or $T_{\mathrm{compliance}}$ actuation depending on
the severity. The MAPE-K loop is a coordination architecture, not a prediction engine.

The Board Projection Holography framework extends this contribution by applying
coordination logic to prospective claims. A board projection is not a statistical
forecast; it is a projection derived from the verified artifact store through
deterministic state transition rules, with every projected state carrying a cryptographic
proof of its provenance. The distinction---proof-backed projection versus statistical
forecast---is the lifecycle project's contribution to the prediction vs.~coordination
debate at the governance layer.

\subsection{Contribution to Process Evidence}
\label{sec:lifecycle-contribution-evidence}

The lifecycle framework's primary contribution to the process evidence corpus is the
formally specified typed evidence lifecycle $\mathcal{S}$ and the six gate criteria
that govern stage advancement. This contribution answers the question: what does it mean
for process evidence to be admissible? The answer is not that the evidence is accurate
or comprehensive in some informal sense; the answer is that it has traversed the typed
state machine from \texttt{Raw} through \texttt{Admitted} under a named witness, and
that every stage transition it has passed through was governed by a formally stated
mathematical predicate.

The Process Asset Taxonomy distinguishes four categories of evidence: Structural Models
(WF-nets, BPMN, POWL), Behavioral Logs (XES, OCEL), Verification Profiles (conformance
reports, alignment scores, soundness receipts), and Cryptographic Receipts (per-action
proofs and decommissioning receipts). The Process Residual Taxonomy identifies what
evidence remains after decommissioning: Structural Residuals (process model fragments
available for reuse), Data Residuals (archived OCEL logs for future discovery runs),
Resource Competence Profiles (operator skill maps derived from the event log), and
Declarative Rules (LTL constraints extracted from the conformance history). Together,
these taxonomies define the full evidence lifecycle from ingestion through archival.

The False-Claim Taxonomy is a negative contribution: it specifies what evidence is not,
by enumerating five structurally inadmissible claim types. This negative specification
is as important as the positive one. Without a formal account of what makes a claim
inadmissible, the admission boundary is vague. The taxonomy gives the boundary a
rigorous form that can be applied in M\&A diligence review.

\subsection{Contribution to Manufacturing Architecture}
\label{sec:lifecycle-contribution-manufacturing}

The lifecycle framework's contribution to the CodeManufactory manufacturing architecture
is the specification of the lifecycle as a manufacturing pipeline in which each stage
is an operator with typed input and typed output. The Design stage manufactures a
\texttt{WfNetConst<SOUNDNESS>} artifact from a schema and declarative constraints. The
Simulation stage manufactures a coverability tree and state-space graph from the WF-net.
The Monitoring stage manufactures \texttt{Metric<Fitness, NUM, DEN>} artifacts from
event streams. The Repair stage manufactures a repaired \texttt{WfNetConst<SOUNDNESS>}
and a \texttt{Receipt}. The Optimization stage manufactures a
\texttt{WfNetConst<SOUNDNESS>} with lower process debt. The Decommissioning stage
manufactures a Cryptographic Decommissioning Receipt $R_d$ and an archived OCEL log.

The Oblivion Protocol---the three-pass ChaCha20 CSPRNG wipe of WASM linear memory
at decommission---is a manufacturing constraint: a decommissioned process model cannot
be re-activated without re-traversing the manufacturing pipeline from Design. This
constraint enforces the immutability doctrine of the CodeManufactory architecture: once
decommissioned, a process model's execution authority is permanently revoked and cannot
be restored without re-manufacturing.

The \texttt{GraduationCandidate} and \texttt{GraduateToWasm4pm} bridge types are the
manufacturing handoff mechanism between the structure manufacturing layer
(\texttt{wasm4pm-compat}) and the execution manufacturing layer (\texttt{wasm4pm}).
An artifact reaches \texttt{GraduationCandidate} status when it has been manufactured
through the full lifecycle specification surface and is ready for execution authority.
The graduation bridge formalizes this handoff as a typed operation, preventing
unauthorized promotion of incomplete artifacts to execution status.

\subsection{Contribution to Capital-Flow Infrastructure}
\label{sec:lifecycle-contribution-capital}

The lifecycle framework's contribution to the capital-flow infrastructure is the
formal grounding of process intelligence as a transactional instrument. The
ACQUISITION.md document establishes that a seller who cannot produce admitted logs and
replayable conformance checks is selling unverifiable narrative, and that the buyer
who cannot replay claims is accepting unquantified risk. This framing treats process
intelligence not as an operational nice-to-have but as a capital-flow instrument: the
presence or absence of a process intelligence package is a direct determinant of
transaction risk premium and, consequently, of valuation.

The Process Debt quantification formula $D_p = \sum_{c \in L}(C_{\mathrm{deviate}}(c)
+ C_{\mathrm{delay}}(c) + C_{\mathrm{overhead}}(c))$ provides the mechanism for
converting process conformance metrics into M\&A valuation adjustments. Process debt
exceeding 15\% of monthly operating cost is defined as the trigger for Debt Actuation
($T_{\mathrm{compliance}}$) and is simultaneously the threshold at which process debt
becomes a material post-close integration liability in M\&A diligence.

The five-layer Process Intelligence Diligence Package---model evidence, log evidence,
conformance evidence, repair evidence, and replay verification---is the capital-flow
artifact that operationalizes this contribution. A seller who can produce all five
layers eliminates the epistemic risk premium. A seller who cannot is at a structural
disadvantage in valuation negotiation. The lifecycle framework is, in this sense, the
formal specification of process intelligence as a capital-flow instrument at the
enterprise transaction layer.
